% ===== §4 =====
\section{Trust}
\label{p27:sec:trust}

\noindent
The first work that relational understanding may do in an intimate bond is the generation of trust, and trust has been theorized with care in several traditions, each of which reaches toward it from its own side. It is worth passing through them, to locate what each reaches and not to survey them, and what each leaves for the account developed here to supply.

A sociological tradition, following Luhmann~\parencite{luhmann1979trust}, treats trust as a mechanism for the reduction of social complexity, a way of absorbing the uncertainty of another's future conduct so that action becomes possible where complete knowledge is not. Trust, on this account, is a wager on the other under conditions of irreducible uncertainty, and it is functional because it permits coordinated life that suspicion would paralyze. A tradition in economics and game theory renders trust in the vocabulary of repeated interaction, reputation, and credible commitment~\parencite{axelrod1984cooperation}, on which trust is the rational expectation that another will not defect, sustained by the shadow of future dealings and the cost of a ruined name. A tradition in the psychology of attachment, following Bowlby and Ainsworth~\parencite{bowlby1969attachment,ainsworth1978patterns}, renders trust as a felt security in the availability of a caregiver, laid down early and carried into the bonds of later life as a working model of whether others can be relied upon.

Each of these renders an aspect of trust, and each leaves an aspect for another to render. The Luhmannian account is precise about the function trust serves and passes over the interior process by which a particular trust is built between two particular people. The game-theoretic account is precise about the conditions under which trust is rational and treats trust as the assessment of a property the other already has, a disposition to cooperate or defect, which the trusting party estimates. The attachment account is precise about the early formation of a capacity for trust and locates its ground in the history of one party and not in the relation between two. Across the three, trust tends to appear as an assessment, by one party, of a property or a probability that belongs to the other or to their circumstances.

Here the two characters of the present account supply what the received theories leave unsupplied, and it is worth marking each in turn. The relational character relocates trust. On the present account trust is not, at its root, an assessment by one party of a property the other already possesses. It is a matter generated between two parties, observable only in the relation between their symbolic worlds, of the same form as any other relational matter. What two people trust, when they trust one another deeply, is something that comes into being in the relation and not a set of estimated dispositions itself and is borne by neither alone. This is why a trust of the deepest kind can be extended before the evidence that a game-theoretic account would require, and why its betrayal wounds in a way the loss of a favorable estimate would not. The trust was not an estimate. It was a shared reality, and its breaking is the breaking of something that was jointly held.

The generative character sets trust in time. On the received accounts trust tends to appear as a state, reached when the conditions are met and maintained while they hold. On the present account trust is generated, deepened, and, where it fails, allowed to wither, through the continued crossing of the two worlds. Each act of understanding, each entry into the other's world that the other recognizes as a true entry, generates a further increment of the shared reality that trust consists in, and a bond in which the partners go on understanding one another is a bond in which trust is continually regenerated and not merely preserved. This sets the account apart from a picture on which trust, once established, is a standing asset that depreciates only under betrayal. Trust that is not regenerated in continued understanding does not hold steady; it thins, as the two worlds drift and the shared reality that trust consisted in is no longer renewed.

The relation of the present account to the received theories is one of location and not of refutation. The Luhmannian function, the game-theoretic conditions, and the attachment history are all real, and the present account does not deny them. It holds that they describe the circumstances, the rationality, and the readiness that surround trust, and that they leave for a relational and generative account the interior matter of what trust is and how it comes to be between two people, which is a shared reality generated in the crossing of their worlds and sustained only as that crossing continues.
